\documentclass[10pt,a4paper]{amsart}
\usepackage[margin=19mm]{geometry}
\usepackage{amsmath,amssymb,mathtools}
\usepackage[scr=rsfs,cal=euler]{mathalfa}
\usepackage{xfrac}
\usepackage[unicode,hidelinks,bookmarks=false]{hyperref}
\newcommand{\ie}{i.e.\ }
\newcommand{\cf}{cf.\ }
\newcommand{\resp}{respectively\ }
\newcommand{\Subsection}{\subsection}
\providecommand{\nobreakdash}{\nobreak\hbox{-}}
\setlength{\emergencystretch}{2em}
\multlinegap0pt
\usepackage{amssymb}
\usepackage{mathtools}
\usepackage{bm}
\usepackage{tikz}
\usepackage[shortlabels]{enumitem}
\newtheorem{proposition}{Proposition}
\newcommand{\define}[1]{\textbf{#1}}
\title{Categorical Perspectives on Chemical Reaction Networks}
\author{Justin Curry, Mauricio Montes}
\date{Source: arXiv:2604.04919v1 (CC BY 4.0)}
\begin{document}
\maketitle

\noindent\emph{Source and licence.} Categorical Perspectives on Chemical Reaction Networks. Version arXiv:2604.04919v1 (CC BY 4.0), distributed by arXiv under the Creative Commons Attribution 4.0 International licence, http://creativecommons.org/licenses/by/4.0/, as recorded in the arXiv licence metadata for this exact version and shown on the abstract page https://arxiv.org/abs/2604.04919v1. Full licence notice: \texttt{licenses/arxiv-2604.04919-CC-BY-4.0.txt}.\par\smallskip

\noindent\textit{Editorial scope.} Counted unit: the article's proposition prop:Schur (source0.tex lines 529--563). The article's complete original proof of it is delivered with it (source0.tex lines 565--643, one \texttt{proof} environment of 79 lines). No mathematics is written, repaired or re-derived: the statement and the proof are the article's own lines, in the article's own notation, and no retained line is substituted or re-flowed.  No further source-local context is needed: every cross-reference of every retained line resolves inside the two delivered spans.  Nothing was omitted from any retained span, so the package carries no omission record.  No retained line contains a citation, so the wrapper carries no bibliography and no bibliography entry.  No retained line contains a figure environment, an image-inclusion command or drawing code, so no image is delivered and the package declares no asset dependency.  Editorial formatting only, in addition: an AMS \texttt{amsart} preamble that restates the article's own declarations and macros used by the retained lines, namely the environments \texttt{proposition} on separate counters, together with the standard packages those lines need.  The retained environments print in this document as Proposition 1 in order of appearance, whereas the article prints the same items with the numbering of its own full document.  The complete native source of the article as distributed in the arXiv e-print for this exact version is bundled unchanged as \texttt{sources/197918/source0.tex}.
\begin{proposition}\label{prop:Schur}
    Let $\Gamma$ be a chemical reaction network and suppose its stoichiometric map in $[\mathbf{A}_2,\mathbf{Vect}]$ admits a block decomposition
    \[
    S =
    \begin{bmatrix}
     a_{11} & a_{12}\\
     a_{21} & a_{22}
    \end{bmatrix}
    :
    C_1^{\mathrm{int}} \oplus C_1^{\mathrm{ext}}
    \longrightarrow
    C_0^{\mathrm{int}} \oplus C_0^{\mathrm{ext}}.
    \]
    
    \begin{enumerate}
        \item If $a_{11}$ is invertible, define the \define{Schur complement} of $a_{11}$ to be
        \[
            \sigma := a_{22} - a_{21}a_{11}^{-1}a_{12}.
        \]
        Then $S$ is isomorphic in $[\mathbf{A}_2,\mathbf{Vect}]$ to the direct sum $a_{11} \oplus \sigma$.
    
        \item More generally, suppose $a_{11}$ may be singular, but the compatibility conditions
        \[
            \ker(a_{11}) \subseteq \ker(a_{21}),
            \qquad
            \operatorname{im}(a_{12}) \subseteq \operatorname{im}(a_{11})
        \]
        hold. Let $a_{11}^+$ be the Moore--Penrose pseudoinverse of $a_{11}$ and define
        \[
            \sigma := a_{22} - a_{21}a_{11}^+a_{12}.
        \]
        Then $\sigma$ is well-defined, independent of the choice of generalized inverse on $\operatorname{im}(a_{12})$, and is the reduced stoichiometric map associated to the elimination of the internal block.
    \end{enumerate}
    In particular, under these hypotheses, the stoichiometric map of the reduced network may be taken to be $S_{\Gamma'} = \sigma$.
\end{proposition}

\begin{proof}
    We begin by rewriting $S_\Gamma$ as a block matrix acting on the "internal" and "external" parts of the chain groups. That is:
    \begin{equation*}
    S = \begin{bmatrix}
        a_{11} & a_{12} \\
        a_{21} & a_{22}
    \end{bmatrix}
    : \overbrace{C_1^{int} \oplus C_1^{ext}}^{C_1(\Gamma)} \to \overbrace{C_0^{int} \oplus C_0^{ext}}^{C_0(\Gamma)}
    \end{equation*}
    Define the inclusions and projections for the left hand side of $S$:
    \begin{align}\label{maps}
        \iota_{\textrm{ext}}^1:C_1^{ext} \to C_1 && \pi_{\textrm{ext}}^1:C_1 \to C_1 ^{\textrm{ext}} \\
        \iota_{\textrm{int}}^1:C_1^{int} \to C_1 && \pi_{\textrm{int}}^1:C_1 \to C_1 ^{\textrm{int}}
    \end{align}
    We define $\iota_{\textrm{ext}}^0, \iota_{\textrm{int}}^0, \pi_{\textrm{ext}}^0, \pi_{\textrm{int}}^0$ in the same way.
    If $a_{11}$ is invertible, then we set $\sigma:=a_{22} -a_{21}a_{11}^{-1}a_{12} :  C_1^{ext} \to  C_0^{ext}$.
    \newline If we consider the following automorphisms:
    \begin{align*}
        L = \begin{bmatrix}
            I & -a_{11}^{-1}a_{12} \\
            0 & I
            \end{bmatrix}:C_1 \to C_1
            &&
        R = \begin{bmatrix}
            I & 0 \\
            -a_{21}a_{11}^{-1} & I
            \end{bmatrix}:C_0 \to C_0   
    \end{align*}
    We notice and set
    \begin{equation*}
        M= RSL =  \begin{bmatrix}
            a_{11} & 0 \\
            0 & \sigma
            \end{bmatrix}
    \end{equation*} 
    Hence, the pairs $(L,R^{-1}):M \Rightarrow S$ and $(L^{-1},R):S \Rightarrow M$ are isomorphisms in the category $[\mathbf{A}_2, \textbf{Vect}]$. From \eqref{maps}, we see that these maps yield morphisms in $[\mathbf{A}_2, \textbf{Vect}]$:
    \begin{align*}
        v_1=(L\iota_{\textrm{int}}^1, R^{-1}\iota_{\textrm{int}}^0): a_{11}\Rightarrow S &&
        v_2=(L\iota_{\textrm{ext}}^1, R^{-1}\iota_{\textrm{ext}}^0): \sigma \Rightarrow S \\
        v_1^{-1}=(\pi_{\textrm{int}}^1L^{-1}, \pi_{\textrm{int}}^0R): S \Rightarrow a_{11} &&
        v_2^{-1}=(\pi_{\textrm{ext}}^1L^{-1}, \pi_{\textrm{ext}}^0R): S \Rightarrow \sigma
    \end{align*}
    These satisfy the property
    \begin{equation*}
        v_1v_1^{-1} + v_2v_2^{-1} = id_S
    \end{equation*}
    Thus, $S$ is isomorphic to the direct sum of $a_{11}$ and $\sigma$ in $[\mathbf{A}_2, \textbf{Vect}]$.

    \noindent In the case where $a_{11}$ is not invertible, we proceed with some additional conditions. Suppose that 
    \[
    \ker a_{11} \subseteq \ker a_{21}, \qquad \textrm{im}(a_{12})\subseteq \textrm{im}(a_{11})
    \]
    The latter condition reveals that for every $x \in C_{1}^{\textrm{ext}}$ there exists a $y \in C_1^{\textrm{int}}$ where 
    \[
    a_{11}y=a_{12}x
    \]
    The former condition implies that $a_{21}y$ is independent of $y$. To see this clearly, for a different $y'$, we see that 
    \[
    a_{11}(y-y') = 0
    \]
    implying $y - y' \in \ker (a_{11}) \subseteq \ker (a_{21})$, thus $a_{21}y = a_{21}y'$. Meaning that 
    \[
    x \longrightarrow a_{22}x - a_{21}y
    \]
    is well defined and defines a linear map.
    Since we had that $\textrm{im}(a_{12}) \subseteq \textrm{im}(a_{11})$. We obtain the following equation:
    \[
    a_{11}a_{11}^{+}a_{12} = a_{12}
    \]
    This means that $a_{12}^{+}a_{12}x$ lifts $a_{12}x$ via $a_{11}$. Hence
    \[
    \sigma (x) = a_{22}x - a_{21}a_{11}^{+}a_{12}x
    \]
    that is,
    \[
    \sigma = a_{22} - a_{21}a_{11}^{+}a_{12}
    \]
    Following the same argument as in the linear case, this formula is independent of the choice of our inverse, provided this agrees on $\textrm{im}(a_{12})$ mod $\ker (a_{11})$. Thus, our reduced stoichiometric matrix is determined by the block decomposition and the compatibility conditions.
\end{proof}

\end{document}
